% !TEX root = ./main.tex

%%%%%%%%%%%%%%%%%%%%%%%%%%%%%%%%%%%%%%%%%%%%%%%%%%%%%%
\section{図表}
図表は文章の一部とはみなされないので，\ref{sec:math}章の数式のように本文中に差し込んではいけません．
図表の配置は
\begin{itemize}
	\item ページの最上段
	\item ページの最下段
	\item 段落が分かれる部分
\end{itemize}
のいずれかに配置します．これはWordで書いても同様です．

%%%%%%%%%%%%%%%%%%%%%%%%%%%%%%%%%%%%%%%%%%%%%%%%%%%%%%
\subsection{図}
%%%%%%%%%%%%%%%%%%%%%%%%%%%%%%%%%%%%%%%%%%%%%%%%%%%%%%
\texttt{pdf} ファイルを取り込むことができます．
\textbf{図~\ref{fig:fig01}} のように，図番号を参照することもできます．
図を貼り付けるときは図の余白に注意（別途説明）．
% /////////////////////////////////////////////////////
\begin{figure}[H]
\centering
	\includegraphics{fig/logo-crop.pdf}
	\caption{舞鶴高専ロゴ}
	\label{fig:fig01}
\end{figure}
% /////////////////////////////////////////////////////

%%%%%%%%%%%%%%%%%%%%%%%%%%%%%%%%%%%%%%%%%%%%%%%%%%%%%%
\subsection{表}
%%%%%%%%%%%%%%%%%%%%%%%%%%%%%%%%%%%%%%%%%%%%%%%%%%%%%%
表の作成は以下の通りです．
書き方のサンプルを2パターン用意しましたのでお好みでどうぞ．
\textbf{表~\ref{table:table02}} の方が，程よい隙間があって見やすい表になります．
また，セル結合の指定なども後者の方がわかりやすく指示できます．
詳細は検索！

% /////////////////////////////////////////////////////
\begin{table}[H]
	\centering
	\caption{旧来からの書き方}\label{table:table01}\vspace*{0.5zh}
\begin{tabular}{|c||c|c|c|c|} \hline
	& 比例感度 \({k}_{P}\) & 比例帯 \({P}_{B}\) & 積分時間 \({T}_{I}\) & 微分時間 \({T}_{D}\) \\ \hline\hline
	P 制御 & 0.5 \({K}_{Pc}\) &	2.0 \({P}_{Bc}\) & --- & --- \\ \hline
	PI 制御 & 0.45 \({K}_{Pc}\) & 2.2 \({P}_{Bc}\) & 0.83 \({T}_{c}\) & --- \\ \hline
	PID 制御 & 0.6 \({K}_{Pc}\) & 1.6 \({P}_{Bc}\) & 0.5 \({T}_{c}\) & 0.125 \({T}_{c}\) \\ \hline
\end{tabular}
\end{table}
% /////////////////////////////////////////////////////
\begin{table}[H]
	\centering
	\caption{最近の書き方}\label{table:table02}\vspace*{0.5zh}
\begin{tblr}{
    hlines, vlines,
    hline{2} = {1}{-}{},hline{2} = {2}{-}{},
	vline{2} = {1}{-}{},vline{2} = {2}{-}{},
    cells={c},
}
& 比例感度 \({k}_{P}\) & 比例帯 \({P}_{B}\) & 積分時間 \({T}_{I}\) & 微分時間 \({T}_{D}\) \\
P 制御 & 0.5 \({K}_{Pc}\) &	2.0 \({P}_{Bc}\) & --- & --- \\
PI 制御 & 0.45 \({K}_{Pc}\) & 2.2 \({P}_{Bc}\) & 0.83 \({T}_{c}\) & --- \\
PID 制御 & 0.6 \({K}_{Pc}\) & 1.6 \({P}_{Bc}\) & 0.5 \({T}_{c}\) & 0.125 \({T}_{c}\) \\
\end{tblr}
\end{table}
